\chapter{Acessibilidade dos Trabalhos Acadêmicos}
\label{chap:acessibilidade}

A acessibilidade permite que o trabalho seja consultado por pessoas com diferentes condições visuais, auditivas, motoras, cognitivas ou tecnológicas. Enquanto a UnDF não publicar requisitos próprios para TCC acessível, as orientações deste capítulo constituem boas práticas editoriais e devem ser aplicadas sempre que a ferramenta permitir.

\boapratica

\section{Documento digital}

O PDF deve conter texto pesquisável, títulos marcados em hierarquia lógica, ordem de leitura coerente, idioma definido, links descritivos e propriedades básicas do documento. Sempre que possível, gere PDF marcado, com estrutura reconhecível por tecnologias assistivas, em vez de digitalizar páginas como imagens.

Use os estilos de título da ferramenta de edição. A aparência visual isolada, como aumentar e colocar o texto em negrito, não cria necessariamente a estrutura que leitores de tela usam para navegar.

\section{Imagens e elementos gráficos}

Toda imagem informativa deve possuir texto alternativo ou descrição equivalente. A descrição comunica finalidade, informação essencial e relação com o argumento, sem repetir mecanicamente a legenda. Imagens decorativas devem ser marcadas como tais quando a ferramenta oferecer essa opção.

Não dependa apenas de cor. Preserve contraste, use rótulos, padrões ou formas adicionais e verifique se o conteúdo permanece compreensível em escala de cinza. Gráficos complexos podem exigir uma síntese textual dos dados mais relevantes.

\section{Tabelas}

Tabelas precisam de cabeçalhos claros e estrutura simples. Evite células mescladas quando prejudicarem a navegação, indique unidades e não use espaços ou tabulações para simular colunas. Tabelas extensas devem ser divididas de forma lógica ou acompanhadas de versão acessível complementar.

\section{Linguagem, fonte e links}

Prefira linguagem clara, frases diretas e termos técnicos explicados na primeira ocorrência. Mantenha fonte legível, tamanho suficiente e espaçamento consistente. Links devem descrever o destino; expressões como ``clique aqui'' não informam o conteúdo fora de contexto.

\section{Conferência}

Antes da entrega, tente selecionar e copiar o texto do PDF, navegue pelos títulos, examine a ordem de leitura, verifique textos alternativos e, quando possível, teste o arquivo com leitor de tela. A validação automática ajuda, mas não substitui a leitura humana.

\section{Slides e pôsteres}

Apresentações devem usar contraste adequado, pouco texto por tela, títulos informativos, ordem de leitura previsível e descrição oral das imagens relevantes. Vídeos precisam de legendas quando utilizados. Pôsteres devem preservar legibilidade à distância e não depender apenas de códigos de cor.

\fluxopendente

A UnDF deve definir requisitos mínimos, ferramentas de teste, responsabilidades e suporte para acessibilidade dos documentos, modelos, apresentações e repositório.
